\noindent\textbf{(a) Resultado.} \[\lambda_n=\left(n+\frac12\right)^2\pi^2,\quad
y_n(x)=C\cos\left((n+\tfrac12)\pi x\right),\quad n=0,1,\ldots,\quad C\ne0.\]

\noindent\textit{Desarrollo.} Multiplicando la ecuación por y e integrando por partes, las condiciones anulan
el término de borde: \(\lambda\int_0^1y^2dx=\int_0^1(y')^2dx\ge0\).
Si \(\lambda=0\), la igualdad obliga a y constante; y(1)=0 la hace nula.
Por tanto \(\lambda=\mu^2>0\). La solución \(A\cos\mu x+B\sin\mu x\)
y la condición inicial dan B=0. Para que A no sea cero se exige
\(\cos\mu=0\), es decir \(\mu=(n+1/2)\pi\), n\(\ge0\).
Cada función indicada satisface la ecuación y ambas condiciones, y la separación de casos excluye otros valores.

\medskip

\noindent\textbf{(b) Resultado.} \(\lambda_0=0\), \(y_0=C\ne0\); para \(n\ge1\),
\(\lambda_n=n^2\pi^2\), \(y_n=C\cos(n\pi x)\), \(C\ne0\).

\noindent\textit{Desarrollo.} El mismo balance \(\lambda\int y^2=\int(y')^2\) descarta autovalores negativos.
Con \(\lambda=0\), las soluciones lineales se reducen por las condiciones a constantes,
y cualquier constante no nula es una autofunción. Si \(\lambda=\mu^2>0\),
y'(0)=0 elimina el seno y y'(1)=0 exige \(\sin\mu=0\).
De ahí \(\mu=n\pi\), n\(\ge1\), y los cosenos de la respuesta.
El modo constante debe conservarse; dividir prematuramente por \(\mu\) lo perdería.
La sustitución directa certifica las funciones propuestas y sus dos condiciones.
